% -----------------------------------------------------------
\section{Iniquity}
% -----------------------------------------------------------
	\label{INIQUITY}
	
% % ----------------------
% \subsection{See also}
% % ----------------------

% ----------------------
\subsection{1828 American Dictionary of the English Language} % *1828
% ----------------------
	\label{INIQUITY-1828}

	\begin{compactitem}
		\item[1] Injustice; unrighteousness; a deviation from rectitude; as the \textit{iniquity} of war; the \textit{iniquity} of the slave trade.
		\item[2] Want of rectitude in principle; as a malicious prosecution originating in the \textit{iniquity} of the author.
		\item[3] A particular deviation from rectitude; a sin or crime; wickedness; any act of injustice.
		\item[4] Original want of holiness or depravity.
	\end{compactitem}

% % ----------------------
% \subsection{Oxford English Dictionary} % *OED
% % ----------------------
	% \label{INIQUITY-OED}

	% \begin{compactitem}
		% \item[]
	% \end{compactitem}

% ----------------------
\subsection{Scriptural Usage} % *SCRIPTURES
% ----------------------
	\label{INIQUITY-SCRIPTURES}

	% % -----
	% \subsubsection{Old Testament} % *OT
	% % -----
		% \label{INIQUITY-OT}
	
		% % --
		% \paragraph{KJV}
		% % --
			% \label{INIQUITY-OT-KJV}
	
			% \begin{tabular}{ll}
				% Total occurrences in the OT & 
			% \end{tabular}
			
			% \begin{compactdesc}
				% \item[]
			% \end{compactdesc}
			
		% % --
		% \paragraph{Hebrew Bible}
		% % --
			% \label{INIQUITY-HEBREW}
		
			% %
			% \subparagraph{\texorpdfstring{\he{sample word}}{}} % paste actual hebrew text here
			% %
				% \label{PAGA} % whatever is in step bible without periods
			
				% \begin{compactdesc}
					% \item[HALOT , PDF ] \textbf{} % Use the 1994 PDF
						% \begin{compactitem}
							% \item[1]
						% \end{compactitem}
					% \item[BDB , PDF ] \textbf{}
						% \begin{compactitem}
							% \item[1]
						% \end{compactitem}
					% \item[DCH PDF ] \textbf{}
						% \begin{compactitem}
							% \item[1]
						% \end{compactitem}
				% \end{compactdesc}
				
				% \begin{compactdesc}
					% \item[Some OT uses] \textbf{}
						% \begin{compactdesc}
							% \item[]
						% \end{compactdesc}
				% \end{compactdesc}
			
		% % --
		% \paragraph{Septuagint}
		% % --
			% \label{INIQUITY-LXX}
		
			% %
			% \subparagraph{\texorpdfstring{\gr{sample word}}{}}
			% %
		
	% -----
	\subsubsection{New Testament} % *NT
	% -----
		\label{INIQUITY-NT}
	
		% --
		\paragraph{KJV}
		% --
			\label{INIQUITY-NT-KJV}
			
	
			\begin{tabular}{ll}
				Total occurrences in the NT & 21
			\end{tabular}
			
			\begin{compactdesc}
				\item[{\hyperref[2THESSALONIANS-2-7]{2 Thessalonians 2:7}}] For the mystery of iniquity doth already work: only he who now letteth will let, until he be taken out of the way.
			\end{compactdesc}
			
		% --
		\paragraph{Greek New Testament} % *GREEK
		% --
			\label{INIQUITY-GREEK}
		
			%
			\subparagraph{\texorpdfstring{\gr{ἀνομία}}{}}
			%
				\label{ANOMIA}
				
				See also \hyperref[HAMARTIA]{\grl{ἁμαρτία}}, \hyperref[PARABASIS]{\grl{παράβασις}}, \hyperref[ADIKIA]{\grl{ἀδικία}},
				
				BDAG and LSJ note that the opposite of \gr{ἀνομία} is \gr{δικαιοσύνη}. \gr{νομία} is listed in LSJ (1179a, PDF 1250), but it's a ``nonce-word'' or a word created for a single occasion and so has just a tiny number of uses. So the lexicons are right to say the proper opposite of \gr{ἀνομία} is \gr{δικαιοσύνη}.
				
				Remember, though, that BDAG also lists \gr{ἀδικία} as an opposite of \gr{δικαιοσύνη}.
			
				\begin{compactdesc}
					\item[BDAG 74b] \textbf{}
						\begin{compactitem}
							\item[1] \textbf{state or condition of being disposed to what is lawless, \textit{lawlessness}}, opposite \gr{δικαιοσύνη}
							\item[2] \textbf{the product of a lawless disposition, \textit{a lawless deed}}
						\end{compactitem}
					\item[CGL 134a, PDF 160] \textbf{}
						\begin{compactitem}
							\item[1] conduct contrary to law, \textbf{lawlessness} (in a region or as a characteristic of an individual or group)...plural, crimes or transgressions \textit{(For the plural meaning, CGL mentions the NT.)}
							\item[2] non-existence or antithesis of law, \textbf{lack} or \textbf{negation of law}
						\end{compactitem}
					\item[LSJ 146b, PDF 219] \textbf{}
						\begin{compactitem}
							\item[1] \textit{lawlessness, lawless conduct}, opposite \gr{δικαιοσύνη}
							\item[2] \textit{the negation of law}, opposite \gr{νόμος}
						\end{compactitem}
				\end{compactdesc}
				
				\begin{tabular}{ll}
					Total occurrences in the NT & 15
				\end{tabular}
				
				\begin{compactdesc}
					\item[All NT uses] \textbf{}
						\begin{compactdesc}
							\item[Matthew 7:23] 
							\item[{\hyperref[MATTHEW-13-41]{Matthew 13:41}}]
							\item[Matthew 23:28]
							\item[Matthew 24:12] 
							\item[Romans 4:7]
							\item[Romans 6:19] 
							\item[{\hyperref[2 Corinthians 6:14]{2 Corinthians 6:14}}] 
							\item[2 Thessalonians 2:3]
							\item[2 Thessalonians 2:7]
							\item[Titus 2:14] 
							\item[Hebrews 1:9] 
							\item[Hebrews 10:17] 
							\item[1 John 3:4]
						\end{compactdesc}
				\end{compactdesc}
		
	% % -----
	% \subsubsection{Book of Mormon} % *BOM
	% % -----
		% \label{INIQUITY-BOM}
	
		% \begin{tabular}{ll}
			% Total occurrences in the BOM & 
		% \end{tabular}
		
		% \begin{compactdesc}
			% \item[]
		% \end{compactdesc}
		
	% % -----
	% \subsubsection{Doctrine and Covenants} % *DC
	% % -----
		% \label{INIQUITY-DC}
	
		% \begin{tabular}{ll}
			% Total occurrences in the DC & 
		% \end{tabular}
		
		% \begin{compactdesc}
			% \item[]
		% \end{compactdesc}
		
	% % -----
	% \subsubsection{Pearl of Great Price} % *PGP
	% % -----
		% \label{INIQUITY-PGP}
	
		% \begin{tabular}{ll}
			% Total occurrences in the PGP & 
		% \end{tabular}
		
		% \begin{compactdesc}
			% \item[]
		% \end{compactdesc}
		
% % ----------------------
% \subsection{Key Phrases} % *KEYPHRASES *PHRASES
% % ----------------------
	% \label{INIQUITY-PHRASES}

	% \begin{compactdesc}
		% \item[] \textbf{\#times} | list
			% % \begin{compactdesc}
				% % \item[Bible anchor verses] \phantom{.}
					% % \begin{compactdesc}
						% % \item[Romans 5:5] \gr{}
					% % \end{compactdesc}
				% % \item[Commentary] ---
			% % \end{compactdesc}
	% \end{compactdesc}
	
% % ----------------------
% \subsection{Bible Dictionaries} % *DICTIONARIES
% % ----------------------
	% \label{INIQUITY-BD}

% % ----------------------
% \subsection{Restoration Usage} % *RESTORATION
% % ----------------------
	% \label{INIQUITY-RESTORATION}

	% % -----
	% \subsubsection{Joseph Smith} % *JS
	% % -----
		% \label{INIQUITY-JS}
	
		% \begin{compactdesc}
			% \item[{\hyperref[KEY]{Date/Source}}]
		% \end{compactdesc}
		
	% % -----
	% \subsubsection{Temple Ordinances}
	% % -----
		% \label{INIQUITY-TEMPLE}
	
		% \begin{compactdesc}
			% \item[{\hyperref[KEY]{Date/Source}}]
		% \end{compactdesc}
	
	% % -----
	% \subsubsection{Brigham Young} % *BY
	% % -----
		% \label{INIQUITY-BY}
	
		% \begin{compactdesc}
			% \item[{\hyperref[KEY]{Date/Source}}]
		% \end{compactdesc}
	
% % ----------------------
% \subsection{Commentary and Studies} % *COMMENTARY *STUDIES
% % ----------------------
	% \label{INIQUITY-COMMENTARY}

	% % -----
	% \subsubsection{Key Points}
	% % -----
		% \label{INIQUITY-KEYS}
	
		% \begin{compactitem}
			% \item
		% \end{compactitem}
		
	% % -----
	% \subsubsection{Sample Study}
	% % -----
		% \label{INIQUITY-study}